\documentclass[12pt]{article}
\usepackage[utf8]{inputenc}
\usepackage[T1]{fontenc}
\usepackage{hyperref}
\title{San Vicente Coatl\'an Zapotec Language Resources for K-12 Curriculum Development}
\author{Compiled by Claude Sonnet 5.5 (claude-sonnet-5-5) for Andr\'as Kornai}
\date{}
\begin{document}
\maketitle

\section{Introduction}
San Vicente Coatl\'an Zapotec (ISO 639-3 zpt; Glottocode sanv1242) is an Otomanguean language of Oaxaca, Mexico. It is also called Southern Ejutla Zapotec or Coatl\'an Zapotec. The Endangered Languages Project rates it as vulnerable, with low certainty. A linguist's comment there says children still acquire it as a first language. Searches found no dictionary, grammar, primer or school textbook for this variety. The resources below are mostly scripture and evangelistic audio.

\section{Primary Educational Resources}
No primers, textbooks or school materials in this language turned up in the searches. Any curriculum would have to be built from scratch, or from the scripture and audio texts listed below.

\section{Linguistic and Reference Materials}
\begin{itemize}
\item Endangered Languages Project entry, with vitality rating and alternate names: \url{https://endangeredlanguages.com/lang/10951}
\end{itemize}
No dictionary or grammar was found.

\section{Religious and Community Resources}
\begin{itemize}
\item Scripture Earth page for the language, with Bible text, a New Testament PDF, video and apps: \url{https://www.scriptureearth.org/00eng.php?iso=zpt}
\end{itemize}
Faith Comes By Hearing lists a New Testament audio recording from Liga B\'iblica Internacional (2012). I did not see a direct URL for it in the search results.

\section{Digital and Audio Resources}
\begin{itemize}
\item Global Recordings Network, audio for the language: \url{https://www.globalrecordings.net/en/language/zpt}
\item Global Recordings Network, alternate page: \url{https://globalrecordings.net/en/language/4063}
\end{itemize}

\section{Curriculum Development Context}
The language appears to be among the healthier Zapotec varieties, but no school use was found. A native-language curriculum would need new orthography, reading materials and teacher-training work with the community. The scripture translation and audio could serve as initial text and audio corpora, with attention to their religious character.

\end{document}